% !TEX root = ../../main.tex
\section{Cấu tạo và nguyên lý cân bằng của quadrotor}
\label{sec:background-principle}

\subsection{Các thành phần chính và cấu hình rôto}
\label{subsec:background-configuration}

Quadrotor sử dụng bốn cụm động cơ--cánh quạt để tạo lực đẩy và điều chỉnh tư thế. Khung máy bay liên kết các cụm động cơ, mang nguồn điện, cảm biến và bộ xử lý điều khiển. Mỗi động cơ nhận công suất qua bộ điều khiển tốc độ động cơ (ESC). Cảm biến cung cấp thông tin về chuyển động của thân; bộ xử lý dùng thông tin này để tính lệnh gửi tới các ESC. Như vậy, khả năng cân bằng phụ thuộc đồng thời vào kết cấu cơ khí, hệ động lực, phép đo và luật điều khiển~\cite{chovancova2014mathematical}.

Để thiết lập mô hình tính toán, xét cấu hình chữ X với bốn rôto bố trí đối xứng quanh trọng tâm. Hệ tọa độ gắn với thân được ký hiệu là $B$: trục $x_B$ hướng về mũi, $y_B$ hướng sang phải và $z_B$ hướng xuống dưới. Góc lăn (roll), góc chúc ngóc (pitch) và góc hướng (yaw) lần lượt mô tả ba thành phần tư thế; ở tư thế gần nằm ngang, chúng tương ứng với chuyển động quay quanh các trục $x_B$, $y_B$ và $z_B$.

Nhìn từ trên xuống, rôto 1 ở phía trước bên trái, rôto 2 ở phía trước bên phải, rôto 3 ở phía sau bên phải và rôto 4 ở phía sau bên trái. Rôto 1 và 3 quay cùng chiều kim đồng hồ (CW); rôto 2 và 4 quay ngược chiều kim đồng hồ (CCW). Đây là quy ước của mô hình, dùng để xác định dấu lực và mômen. Hình học bố trí rôto, chiều quay và các trục thân được minh họa trong \fref{fig:background-layout}.
% Cần bổ sung: đối chiếu cấu hình khung, số thứ tự và chiều quay động cơ với mô hình lắp ráp thực tế.

\begin{figure}[H]
    \centering
    \resizebox{\linewidth}{!}{\begin{tikzpicture}[
    x=1cm,y=1cm,
    >=latex,
    line width=0.8pt,
    font=\fontsize{11}{13}\selectfont,
    every node/.style={align=center}
]
    % Schematic convention for this figure; viewed from above.
    \draw[gray!55,line width=8pt,line cap=round] (-2,-2) -- (2,2);
    \draw[gray!55,line width=8pt,line cap=round] (-2,2) -- (2,-2);
    \draw[fill=gray!10,rounded corners=2pt] (-0.48,-0.48) rectangle (0.48,0.48);

    % Opposite rotors rotate in the same direction.
    \foreach \x/\y/\n in {-2/2/1,2/-2/3}{
        \begin{scope}[shift={(\x,\y)}]
            \draw[fill=white] (0,0) circle (0.73);
            \draw[->,line width=1pt] (130:0.87) arc (130:-130:0.87);
            \node at (0,0.19) {Rôto \n};
            \node at (0,-0.22) {CW};
        \end{scope}
    }
    \foreach \x/\y/\n in {2/2/2,-2/-2/4}{
        \begin{scope}[shift={(\x,\y)}]
            \draw[fill=white] (0,0) circle (0.73);
            \draw[->,line width=1pt] (50:0.87) arc (50:310:0.87);
            \node at (0,0.19) {Rôto \n};
            \node at (0,-0.22) {CCW};
        \end{scope}
    }

    % Body frame: x points forward, y points right, z points into the page.
    \draw[->] (0,0) -- (0,3.15) node[above] {$x_B$ -- hướng mũi};
    \draw[->] (0,0) -- (3.6,0) node[right] {$y_B$};
    \draw[fill=white] (0,0) circle (0.13);
    \draw (-0.085,-0.085) -- (0.085,0.085);
    \draw (-0.085,0.085) -- (0.085,-0.085);
    \node[anchor=north west] at (0.52,-0.32) {$z_B$};

    \node[anchor=west,align=left] at (3.15,2.0)
        {Roll: quay quanh $x_B$\\Pitch: quay quanh $y_B$\\Yaw: quay quanh $z_B$};
    \node[anchor=west,align=left] at (3.15,-1.7)
        {$\otimes\;z_B$: hướng vào trang\\CW: cùng chiều kim đồng hồ\\CCW: ngược chiều kim đồng hồ};
\end{tikzpicture}
}
    \caption{Quy ước cấu hình quadrotor chữ X, chiều quay rôto và hệ trục tọa độ thân khi nhìn từ trên xuống.}
    \label{fig:background-layout}
\end{figure}

\subsection{Nguyên lý tạo lực đẩy và mômen cân bằng}

Lực đẩy của mỗi rôto có phương gần trùng với trục $z_B$ và chiều hướng lên trên thân, tức ngược chiều dương $z_B$. Thay đổi đồng thời lực đẩy của bốn rôto làm thay đổi tổng lực đẩy. Khi thân nghiêng, hướng tổng lực đẩy cũng nghiêng theo, tạo thêm thành phần lực theo phương ngang. Vì vậy, điều chỉnh tư thế đồng thời làm thay đổi các thành phần gia tốc tịnh tiến~\cite{bouabdallah2004design}.

Chênh lệch lực đẩy giữa hai phía trái và phải tạo mômen roll; chênh lệch giữa phía trước và phía sau tạo mômen pitch. Với hệ trục đã quy ước, tăng lực đẩy phía trái so với phía phải tạo mômen roll dương, còn tăng lực đẩy phía trước so với phía sau tạo mômen pitch dương. Trong cấu hình chữ X, mỗi rôto đều có cánh tay đòn đối với cả hai trục nên việc hiệu chỉnh tư thế sử dụng sự phối hợp của bốn động cơ.

Mômen yaw hình thành từ chênh lệch mômen phản lực giữa hai nhóm rôto quay ngược chiều. Ở chế độ quay gần xác lập, mômen truyền từ động cơ lên thân có chiều ngược với chiều dẫn động rôto. Do đó, rôto 1 và 3 tạo mômen yaw âm theo $z_B$; rôto 2 và 4 tạo mômen yaw dương. Khi các mômen phản lực của hai nhóm cân bằng nhau, tổng mômen yaw bằng không.

Cân bằng tư thế là bài toán duy trì hoặc đưa các góc quay về giá trị đặt bằng cách điều chỉnh mômen tác dụng lên thân. Điều kiện này khác với cân bằng lực khi bay treo. Với máy bay tự do ở tư thế nằm ngang, bỏ qua lực ngoài khác, tổng lực đẩy bằng trọng lực mới cho gia tốc tịnh tiến bằng không. Vì vậy, giữ roll và pitch gần không chưa đủ để duy trì độ cao hoặc vị trí. Đối với bộ điều khiển cân bằng, tổng lực đẩy và các mômen là những đại lượng đầu vào có vai trò khác nhau.

\subsection{Quan hệ giữa tốc độ rôto, lực đẩy và mômen}

Trong mô hình khí động học tĩnh giản lược, lực đẩy và độ lớn mômen cản khí động của mỗi rôto tỉ lệ với bình phương tốc độ quay~\cite{chovancova2014mathematical}:
\begin{equation}
    f_i=c_T\Omega_i^2,\qquad Q_i=c_Q\Omega_i^2,
    \qquad i=1,2,3,4.
    \label{eq:background-rotor-force}
\end{equation}
Ở đây, $\Omega_i$ là tốc độ góc rôto, tính bằng rad/s; $f_i$ là độ lớn lực đẩy, tính bằng N; $Q_i$ là độ lớn mômen cản, tính bằng N$\cdot$m. Các hệ số $c_T>0$ và $c_Q>0$ phụ thuộc vào cụm động cơ--cánh quạt và điều kiện làm việc. Mô hình sử dụng chung hai hệ số cho bốn rôto khi xem các cụm động lực có đặc tính tương đương. Trong miền vận hành đang xét, các hệ số được coi là không đổi.
% Cần bổ sung: thông số hoặc dữ liệu đo để xác định c_T, c_Q cho cụm động cơ--cánh quạt sử dụng.

Tổng lực đẩy có độ lớn:
\begin{equation}
    T=\sum_{i=1}^{4}f_i=c_T\sum_{i=1}^{4}\Omega_i^2.
    \label{eq:background-total-thrust}
\end{equation}
Biểu thức này cho biết ảnh hưởng của tốc độ rôto đến lực đẩy. Đối với chuyển động quay, mômen còn phụ thuộc vào vị trí đặt lực và chiều quay từng rôto. Do đó, để xác định tư thế biến đổi như thế nào khi các lực đẩy thay đổi, cần mô tả chuyển động của thân bằng các quan hệ động học và động lực học vật rắn.
